\section{Introduction}
\label{sec:intro}



The art of computer programming is naturally an interactive process. When a human programmer writes code, she relies on several iterations of a `write-execute-test' loop in order to iteratively refine solutions, plan changes, test sub-modules, and solve ambiguities by checking execution behavior. While this is reminiscent of other human endeavors like writing, code compilation and execution produce exact results that provide a deterministic form of feedback to make the refinement process more straightforward. Depending on the observed results, programmers perform various levels of debugging and rewriting, and continue the process until their code satisfies the requirements.

There has been increasing interest in recent years around the development of models that can automatically generate code given a specification in natural language~\cite{feng2020codebert, wang2021codet5, clement2020pymt5, alphacode, le2022coderl}. Powered by large-scale pre-training over thousands of codebases~\cite{Agashe2019JuICeAL, husain2020codesearchnet, gao2020pile}, these models have shown solid performance on static benchmarks like HumanEval~\cite{chen2021evaluating}, APPS \cite{hendrycksapps2021}, MBPP \cite{austin2021program}, CodeXGLUE \cite{lu2021codexglue}. However, generating code in a static, sequence-to-sequence or auto-regressive fashion has several drawbacks: 1) simple errors (even typos) can propagate and there is no chance for recovery or revision, 2) there is a disconnect between the code generation process and its downstream execution on the desired software and hardware environment, and 3) there is little room for human intervention or collaboration in the code generation process. 

\begin{figure}
    \centering
    \includegraphics[width=1.0\textwidth]{figures/assets/Preview.pdf}
    \caption{Overview of \benchmark. Setting up an interactive code environment with \benchmark{} requires a Dockerfile, dataset, reward function definition, and a small amount of subclass implementation. The interactive loop between agent and environment closely mirrors real world software development processes. While \benchmark{} task performance is generally quantified as a binary 0/1 completion score, \benchmark{} allows for the design of more complex evaluation criteria that can incorporate execution output and the effects of interaction on the state space.}
    \label{fig:preview}
    \vspace{-1.5em}
\end{figure}

Recently, some works have proposed the use of execution feedback or interaction~\cite{wang2023interactive} to benefit code generation models~\cite{Lai2022DS1000, huang-etal-2022-execution, wang2023executionbased, hendrycksapps2021}. However, these papers consider their own individual setup and are difficult to compare with one other due to the use of different compilers, execution environments, feedback signals, and assumptions on the interactive process such as human participation to create task descriptions or provide natural language feedback. This makes it difficult to compare existing methods for code generation and to clearly understand the benefits of interactive generation. 

To address these issues, we propose InterCode, the first standard coding benchmark designed natively with an interactive execution environment. Closely mimicking the human decision-making process, \benchmark{} allows a coding agent to interactively receive feedback from compilers/interpreters that execute its code, and to submit further refinements. We design \benchmark{} to be like a standard reinforcement learning (RL) environment that requires minimal human intervention and one in which generated code is treated as actions, which are executed to reveal observations. Our framework is (1) language and platform agnostic and can easily be used for new coding problems, (2) uses self-contained Docker environments to provide safe execution, and (3) compatible out-of-the-box with traditional seq2seq generation methods, while also enabling and empowering the development of new interactive techniques. 

We demonstrate the power of the framework by implementing SQL and bash tasks within \benchmark{}, building on pre-existing static datasets~\cite{zhong2017seq2sql, lin-etal-2018-nl2bash}. We perform experiments across diverse models and prompting methods, including ReAct~\cite{yao2023react} and Plan \& Solve~\cite{wang2023planandsolve}. Our findings concretely showcase the benefits of interaction towards solving coding tasks, discuss the distribution of distinct code understanding challenges across different task settings, and explore the ease with which new tasks and datasets can be defined using \benchmark{}.

To summarize, our paper makes the following contributions:
\begin{itemize}%
\item We develop \benchmark{}, a new, universal framework for interactive code generation, which provides ease of use, extensibility, and safety.
\item Using \benchmark{}, we perform a comprehensive evaluation of state-of-the-art models and identify several avenues for improvements.
\item We release our framework as a new benchmark along with useful empirical tools to customize any new static code datasets into interactive tasks.
\end{itemize}












